\newcommand{\srEq}{\mathsf{Eq}}
\newcommand{\srInc}{\mathsf{Inc}}
\newcommand{\srCon}{\operatorname{Con}}
\newcommand{\srcl}[1]{\langle #1\rangle}
\newcommand{\srrho}{\rho}
\newtheorem{srtheorem}{Theorem}
\newtheorem{srlemma}[srtheorem]{Lemma}

\begin{abstract}
The source44 nine-function closure kernels FC15$^*$ and FC16$^*$ have
respectively 29 and 30 universal clauses after treating equality as a
predicate and adding the published congruence completion. Seven function
congruence clauses can be removed while preserving consistency, provably
in ZFC. The resulting ground search retains function congruence only for
assignment evaluation $V$ and assignment update $C$. Class constructors are reconstructed by
their incidence profiles; witness functions are reconstructed after the
equality quotient. The resulting exact payloads have 22 clauses, 58
literals and 1,195 template bits, or 23 clauses, 62 literals and 1,402
bits for the extensional control. Two further exact redundancies, for the
branch guard and explicit reflexivity, give final payloads of 20 clauses,
53 literals and 1,134 bits, or 21 clauses, 57 literals and 1,341 bits.
The reduction is a theorem about a
finite formalism and its refutation-search specification. No binary
Turing machine or state saving is inferred from these counts.
\end{abstract}

\section{Exact target and source dependency}

The starting point is the recommended, guard-fused, derived-member version
of source44, not its twelve-function explanatory presentation. Its fixed
signature has one constant $d$, the nine functions
\[
V^2,C^3,N^2,A^2,Q^2,W^3,Y^3,Z^3,H^2,
\]
and the binary predicates $E,\srInc,G$. Native FC15$^*$ uses genuine
logical equality. The ground CNF representation instead uses the binary
predicate $\srEq$. Superscripts indicate arity. All variables in every
displayed clause are universally quantified over one common domain;
colliding substitutions are included.

Set
\begin{align*}
q&:=N(d,d),\\
R(b,t,x,y)&:=\srInc(C(C(b,d,x),q,y),t),\\
D^*(x)&:=W(C(d,q,x),d,A(d,q)).
\end{align*}
The fifteen native rows, with equality written as $\srEq$ for the ground
representation, are exactly the following:
\begin{align}
&\srEq(V(C(s,j,z),j),z),\tag{01}\\
&\srEq(k,j)\lor\srEq(V(C(s,j,z),k),V(s,k)),\tag{02}\\
&\srInc(s,N(a,b))\lor\srInc(s,a),\tag{03}\\
&\srInc(s,N(a,b))\lor\srInc(s,b),\tag{04}\\
&\neg\srInc(s,N(a,b))\lor\neg\srInc(s,a)\lor\neg\srInc(s,b),\tag{05}\\
&\neg\srInc(s,A(j,k))\lor E(V(s,j),V(s,k)),\tag{06}\\
&\srInc(s,A(j,k))\lor\neg E(V(s,j),V(s,k)),\tag{07}\\
&\neg\srInc(C(s,j,z),a)\lor\srInc(s,Q(j,a)),\tag{08}\\
&\neg\srInc(s,Q(j,a))\lor\srInc(C(s,j,W(s,j,a)),a),\tag{09}\\
&G(b,t)\lor R(b,t,H(b,t),Y(b,t,u)),\tag{10}\\
&G(b,t)\lor E(Z(b,t,u),Y(b,t,u)),\tag{11}\\
&G(b,t)\lor\neg E(Z(b,t,u),u),\tag{12}\\
&\neg G(b,t)\lor\neg E(x,H(b,t))\lor\neg R(b,t,x,y)
       \lor E(y,H(b,t)),\tag{13}\\
&\neg G(b,t)\lor E(D^*(x),x)\lor\neg R(b,t,x,y)
       \lor E(y,H(b,t)),\tag{14}\\
&\neg G(b,t)\lor\neg E(y,H(b,t))\lor R(b,t,Y(b,t,y),y).\tag{15}
\end{align}
These are universal CNF clauses, not Horn rules: rows 02--04, 10--11 and
14 have more than one positive literal. In particular, satisfiability
and propositional compactness, rather than least Horn closure, are the
correct semantic tools.

For a list $\mathcal F$ of functions, let $\mathrm{EC}(\mathcal F)$ contain
the following equality completion:
\begin{align*}
&\srEq(x,x),\\
&\neg\srEq(x,y)\lor\neg\srEq(x,z)\lor\srEq(y,z),\\
&\bigvee_{i<r}\neg\srEq(x_i,y_i)
  \lor\srEq(F(x_0,\ldots,x_{r-1}),F(y_0,\ldots,y_{r-1}))
  &&(F^r\in\mathcal F),\\
&\neg\srEq(x_0,y_0)\lor\neg\srEq(x_1,y_1)
  \lor\neg P(x_0,x_1)\lor P(y_0,y_1)
  &&(P\in\{E,\srInc,G\}).
\end{align*}
The two equality laws give equivalence, and the last three rows make the
predicates constant on equivalence classes. Our target $U_{15}^{VC}$
consists of rows 01--15 and $\mathrm{EC}(\{V,C\})$. Thus it has seven
completion rows: two equality laws, two function rows, and three predicate
rows. No native row or predicate-congruence row is deleted.

The source dependence is deliberately modular. Source44 supplies the
ordinary ZFC-internalizable calibration
\[
\srCon(\mathrm{FC15}^*)\ \longleftrightarrow\ \srCon(\mathrm{ZFC}),
\qquad
\srCon(\mathrm{FC16}^*)\ \longleftrightarrow\ \srCon(\mathrm{ZFC}).
\]
The FC15$^*$ lower direction names the earlier full raw-EPSC hereditary
interpretation; FC16$^*$ has the separate extensional lower route. The new
proof below is a direct model reconstruction between the finite theories.
It does not reprove or silently strengthen those inherited calibrations.

\section{A quotient with reconstructed operations}

Let $\mathcal T$ be the countable set of ground terms in the nine-function
signature, with the source's seven unused unary padding constructors if
using its total radix16 codec. Terms are ordered by their natural-number
codes, with $d$ of code zero. Padding receives no congruence requirements.
Assume a Boolean valuation $v$ satisfies every ground instance of
$U_{15}^{VC}$. This is obtained in ZFC from consistency by completeness,
or from finite ground satisfiability by propositional compactness.

Define $t\sim u$ iff $v(\srEq(t,u))=1$, and let
$\mathcal D=\mathcal T/{\sim}$. Write $\srcl t$ for the class of $t$.
Each class is nonempty and consists of natural-number codes, so it has a
least-coded representative $\srrho(a)$. This explicitly defined selector
requires no arbitrary choice of an uncountable family. In particular,
$\srrho(\srcl d)=d$.

Predicate congruence makes
\[
P^{\mathcal D}(\srcl t,\srcl u)=v(P(t,u))
\quad (P=E,\srInc,G)
\]
well-defined. Keep the two assignment operations by ordinary quotient:
\[
V^{\mathcal D}(\srcl t,\srcl u)=\srcl{V(t,u)},\qquad
C^{\mathcal D}(\srcl t,\srcl u,\srcl w)=\srcl{C(t,u,w)}.
\]
Their retained congruence rows are exactly what makes these definitions
independent of the chosen terms.

For every other function $F$, define a \emph{new} operation by
\begin{equation}\label{sr-reconstruct}
F^{\mathcal D}(a_1,\ldots,a_r)
  :=\srcl{F(\srrho(a_1),\ldots,\srrho(a_r))}.
\end{equation}
Set $d^{\mathcal D}=\srcl d$. Definition
\eqref{sr-reconstruct} is meaningful for $N,A,Q,W,Y,Z,H$ without any
congruence of the corresponding raw term constructors. It is not a
quotient homomorphism for those constructors. In general one must \emph{not}
assert
\[
F^{\mathcal D}(\srcl{t_1},\ldots,\srcl{t_r})
   =\srcl{F(t_1,\ldots,t_r)}
\]
when arbitrary $t_i$ replace the least representatives. The proof proceeds
by the defining contracts of the operations and by the placement of their
witness occurrences, instead of making this false inference.

\begin{srlemma}[Core reconstruction]\label{sr-core}
The reconstructed structure satisfies native rows 01--09 with genuine
equality. It satisfies
\begin{align*}
\srInc(s,N(a,b))&\longleftrightarrow
                  \neg(\srInc(s,a)\land\srInc(s,b)),\\
\srInc(s,A(j,k))&\longleftrightarrow E(V(s,j),V(s,k)),\\
\srInc(s,Q(j,a))&\longleftrightarrow
                  \exists z\,\srInc(C(s,j,z),a),
\end{align*}
where every symbol in these three displayed contracts is interpreted in
$\mathcal D$. Moreover $q^{\mathcal D}=\srcl{N(d,d)}$.
\end{srlemma}
\begin{proof}
For rows 01--02, choose representatives of the universal variables and use
only quotient-preserved $V,C$. Ground equality then becomes actual equality
of classes. For rows 03--05 fix the least representatives of $a,b$ and
instantiate the three rows at any representative of $s$. Their Boolean
combination gives precisely the NAND profile of the reconstructed
$N(a,b)$. For rows 06--07 fix the representatives of $j,k$; the raw $V$
applications have the required quotient values, proving the atom profile.

For $Q(j,a)$ fix $j_0=\srrho(j)$ and $a_0=\srrho(a)$ and, for a given
$s$, set $s_0=\srrho(s)$. If $\srInc(s,Q(j,a))$ holds, raw row 09 at
$(s_0,j_0,a_0)$ supplies the ground witness $W(s_0,j_0,a_0)$. Its class is
exactly the reconstructed $W(s,j,a)$ by \eqref{sr-reconstruct}; preservation
of $C$ and predicate congruence give native row 09. Conversely each $z$ has
a ground representative, and raw row 08 at that representative gives the
reverse direction and native row 08. All these facts hold for every $s$,
not just the least representative as a syntactic term. Finally both
arguments used to reconstruct $N(d^{\mathcal D},d^{\mathcal D})$ have least
representative $d$, proving the assertion about $q$.
\end{proof}

Two core consequences will be used at the only problematic nested terms.
First, NAND at $(d,d)$ gives $d^{\mathcal D}\ne q^{\mathcal D}$: equality
would identify a proposition with its negation. Thus
\[
P(x,y):=C(C(d,d,x),q,y)
\]
has projections $V(P(x,y),d)=x$ and $V(P(x,y),q)=y$ and is injective in
both coordinates jointly. Put $a_0=P(q,d)$ and $a_{n+1}=P(d,a_n)$.
The map $P(d,-)$ is injective and $a_0$ is outside its range, so its finite
iterates are pairwise different. In particular there is a selector
$j\notin\{d,q\}$; more generally there are as many distinct selectors as
any fixed finite formula needs. This is a finite equational argument, not
an assumption that one fixed named selector avoids $d,q$ in every model.

Second, in \emph{every} structure satisfying the reconstructed core,
\begin{equation}\label{sr-member}
(\exists z\ E(z,x))\ \longrightarrow\ E(D^*(x),x).
\end{equation}
To check this, let $s=C(d,q,x)$. Since $d\ne q$, updating its $d$ coordinate
to a member $z$ of $x$ gives values $z,x$ at $d,q$. The atom and existential
contracts give $\srInc(s,Q(d,A(d,q)))$. Native row 09 with these actual
reconstructed arguments supplies $W(s,d,A(d,q))=D^*(x)$, and the atom
contract identifies it as a member of $x$.

\begin{srlemma}[Joint reconstruction of the closure rows]\label{sr-closure}
The reconstructed structure satisfies native rows 10--15.
\end{srlemma}
\begin{proof}
Fix $b,t\in\mathcal D$, once and for all, and write
$b_0=\srrho(b),t_0=\srrho(t)$ and $h_0=H(b_0,t_0)$. Thus the reconstructed
$H(b,t)$ is $\srcl{h_0}$ in every one of the six rows. The guard is the
single fixed value $v(G(b_0,t_0))$. Choosing representatives separately for
the rows could lose this shared witness key; the present construction does
not do that.

All occurrences of $R$ preserve their meanings under the quotient, because
their only proper operations are $C$ and the fixed ground abbreviation
$q=N(d,d)$, whose class is exactly the reconstructed $q$. For example, for
arbitrary ground $x_0,y_0$,
\[
v(R(b_0,t_0,x_0,y_0))
 =R^{\mathcal D}(b,t,\srcl{x_0},\srcl{y_0}).
\]
The equality holds also when $x_0$ or $y_0$ is a raw witness application;
that application is used as a ground argument, not assumed to commute with
evaluation in the reconstructed language.

If $G(b,t)$ is false, fix any $u$ and its representative $u_0$. The raw
instances of rows 10--12 share $h_0,Y(b_0,t_0,u_0),Z(b_0,t_0,u_0)$. Their
classes are exactly the reconstructed $H(b,t),Y(b,t,u),Z(b,t,u)$.
Preservation of $R,E$ gives all three native rows. Rows 13--15 are already
true through their guard.

If $G(b,t)$ is true, rows 10--12 are guarded. For arbitrary $x,y$ choose
$x_0=\srrho(x),y_0=\srrho(y)$. Raw row 13, with the same $b_0,t_0$, gives
native row 13 because its witness is the same $h_0$. Raw row 15 gives
native row 15 with reconstructed input $Y(b,t,y)$, since it uses exactly
$Y(b_0,t_0,y_0)$.

For row 14 distinguish two cases. If $x$ has a member in $\mathcal D$,
the \emph{reconstructed} term $D^*(x)$ is a member by
\eqref{sr-member}, so its positive literal makes the row true. If $x$ has
no member, consider instead the literal in the \emph{raw} ground row 14:
\[
E\bigl(W(C(d,N(d,d),x_0),d,A(d,N(d,d))),x_0\bigr).
\]
Its value is false because the first endpoint has some class in
$\mathcal D$ and no class belongs to $x$. The other literals are the
preserved guard, $R$, and $E(y,H(b,t))$, so the raw row gives the desired
closure and hence native row 14. There is no assertion that this raw
$D^*$ term has the same class as the reconstructed evaluation of $D^*(x)$.
\end{proof}

\begin{srtheorem}[Finite-kernel equisatisfiability]\label{sr-main}
The theories $U_{15}^{VC}$ and native FC15$^*$ have models simultaneously,
by constructions provable in ZFC. In particular, ZFC proves
\[
\srCon(U_{15}^{VC})\ \longleftrightarrow\ \srCon(\mathrm{FC15}^*).
\]
Consequently, using the named source44 calibration, ZFC proves
$\srCon(U_{15}^{VC})\leftrightarrow\srCon(\mathrm{ZFC})$.
\end{srtheorem}
\begin{proof}
A native FC15$^*$ model interprets $\srEq$ as genuine equality and satisfies
every retained completion row, so it gives a model of $U_{15}^{VC}$. For
the other direction obtain an all-ground valuation, form the quotient
above, and use Lemmas~\ref{sr-core} and~\ref{sr-closure} to obtain a full
native FC15$^*$ model. Completeness in each direction converts model
existence to syntactic consistency. Ground terms, valuation, quotient,
least representatives and finite syntax are set-coded objects in ZFC.
No claim of a serialized proof-assistant proof is made.
\end{proof}

\section{The 23-clause extensional control}

Define expanded terms
\[
\operatorname{Xor}(u,v):=N(N(u,N(u,v)),N(v,N(u,v))),\qquad
T_j:=\operatorname{Xor}(A(j,d),A(j,q)).
\]
FC16$^*$ adds to the fifteen native rows the following four-literal row:
\begin{equation}\label{sr-ext-row}
\srEq(j,d)\lor\srEq(j,q)\lor\srEq(a,b)
 \lor\srInc(C(C(d,d,a),q,b),Q(j,T_j)).
\end{equation}
Let $U_{16}^{VC}$ add that same fully expanded raw clause to
$U_{15}^{VC}$. It introduces no new function.

\begin{srlemma}[Raw profiles recover extensionality]\label{sr-ext}
The reconstructed quotient of every all-ground model of $U_{16}^{VC}$ has
extensional $E$: distinct objects have a distinguishing incoming member.
\end{srlemma}
\begin{proof}
Choose any $a\ne b$ in the quotient and a selector $j\notin\{d,q\}$,
which exists by the reconstructed core. Let $a_0,b_0,j_0$ be ground
representatives. Put
\[
s_0=C(C(d,d,a_0),q,b_0),\qquad
t_0=\operatorname{Xor}(A(j_0,d),A(j_0,q)),
\]
where these are \emph{raw ground terms}, with $q$ expanded as $N(d,d)$.
The three equality literals of the corresponding ground instance of
\eqref{sr-ext-row} are false, so $v(\srInc(s_0,Q(j_0,t_0)))=1$.

Raw rows 08--09, with these exact ground arguments, say
\[
v(\srInc(s_0,Q(j_0,t_0)))=1
 \quad\Longleftrightarrow\quad
\exists z_0\in\mathcal T\quad
 v(\srInc(C(s_0,j_0,z_0),t_0))=1.
\]
For a ground $z_0$, repeated raw instances of rows 03--05 evaluate the
displayed NAND tree as exclusive OR of its two atom incidences. Raw rows
06--07 evaluate these atom incidences as
\[
E^{\mathcal D}(\srcl{z_0},a),\qquad
E^{\mathcal D}(\srcl{z_0},b).
\]
Indeed $V,C$ are preserved, and updating the assignment $s_0$ at the fresh
selector $j_0$ sets its value there to $\srcl{z_0}$ while retaining $a,b$
at $d,q$. Therefore the true final literal supplies a distinguishing
member. This argument is an induction on the \emph{raw incidence profiles};
it does not identify $\srcl{t_0}$ with the value of the reconstructed
$T_j$, or $\srcl{Q(j_0,t_0)}$ with the corresponding reconstructed term.
\end{proof}

Over the reconstructed core, extensionality now implies the actual native
row \eqref{sr-ext-row}: if none of its equality guards holds, choose a
distinguishing member and apply the reconstructed atom, NAND and existential
contracts. Thus the reconstructed structure is a native FC16$^*$ model,
not merely a structure satisfying an unexpanded informal schema. The
reverse model implication again uses genuine equality. We obtain
\[
\mathrm{ZFC}\vdash
\srCon(U_{16}^{VC})\leftrightarrow\srCon(\mathrm{FC16}^*)
\leftrightarrow\srCon(\mathrm{ZFC}),
\]
where the last equivalence has the same inherited source44 status as before.

\section{Two further completion rows can be removed}

\subsection{Reconstructing the branch predicate}

The guard $G$ need not itself descend through the equality quotient.
Delete its four-literal predicate-congruence row, retain the congruence
rows for $E$ and $\srInc$, and define instead
\[
G^{\mathcal D}(b,t):=v(G(\srrho(b),\srrho(t))).
\]
Every occurrence of $G$ in the native theory is the variable-only shared
key $(b,t)$ of rows 10--15. Lemma~\ref{sr-closure} already fixes exactly
these two representatives jointly before considering the branch and its
$H,Y,Z$ witnesses. Therefore its proof applies verbatim with the new
guard definition. The core, raw $D^*$ emptiness step and FC16$^*$
extensionality argument never require transport of $G$ between arbitrary
raw terms. The original native model still gives the reverse direction.
This removes one row, four literals, eight term nodes and exactly 47
template bits: the intermediate targets have 21/54/1,148 and
22/58/1,355 clauses/literals/bits.

This observation uses the exact occurrence pattern of the guard, not a
general license to omit predicate congruence. In particular $E$ and
$\srInc$ remain quotient predicates in the proof.

\subsection{Reflexivity is already a consequence of native row 01}

For an arbitrary ground term $z$, set $p_z=V(C(d,d,z),d)$. The instance
of native row 01 with $s=j=d$ gives $\srEq(p_z,z)$. The retained
left-Euclidean row at $(x,y,z)=(p_z,z,z)$ is
\[
\neg\srEq(p_z,z)\lor\neg\srEq(p_z,z)\lor\srEq(z,z).
\]
Together these derive $\srEq(z,z)$. Thus all-ground reflexivity is
available without its explicit one-literal template. Symmetry and
transitivity then follow from reflexivity and left Euclidean equality
as before. This is an actual two-premise derivation, not an assumption
that the equality predicate is reflexive by its name.

Let $U_{15}^{\circ}$ contain the fifteen native clauses and exactly five
completion clauses: left Euclidean equality, function congruence for
$V,C$, and predicate congruence for $E,\srInc$. Let $U_{16}^{\circ}$
also include the unchanged native extensionality clause. Derive global
reflexivity first; form the quotient and reconstruct functions and guard
as above. The proofs yield
\[
\begin{split}
U_{15}^{\circ}\text{ has a model}
  &\quad\Longleftrightarrow\quad\mathrm{FC15}^*\text{ has a model},\\
U_{16}^{\circ}\text{ has a model}
  &\quad\Longleftrightarrow\quad\mathrm{FC16}^*\text{ has a model}.
\end{split}
\]
Both constructions are internalizable in ZFC. Applying the named native
calibrations gives
\[
\mathrm{ZFC}\vdash
\srCon(U_{15}^{\circ})\leftrightarrow\srCon(\mathrm{ZFC}),
\qquad
\mathrm{ZFC}\vdash
\srCon(U_{16}^{\circ})\leftrightarrow\srCon(\mathrm{ZFC}).
\]
Removing the reflexivity template saves one clause, one literal, two
term nodes and 14 more bits. The final targets are 20 clauses, 53
literals and 1,134 bits, or 21 clauses, 57 literals and 1,341 bits.
These are the smallest completions proved in this contribution; no
minimality theorem is claimed.

\section{Exact encoding comparison}

The companion generator uses byte-exact copies of source44's
\texttt{ground\_kernel.py} and \texttt{coded\_extensional.py}. It calls the
existing clause generator unchanged and varies only the function list
passed to its equality-completion routine. The original prefix template
format uses a three-bit clause length, a sign bit and a two-bit predicate
tag per literal, four-bit prefix term opcodes, and a terminating zero
clause length. Six variable opcodes, one constant and nine function
opcodes occupy all sixteen term codes. The same nine-function ground
enumeration remains present in every variant.

\begin{center}\small
\begin{tabular}{llrrrr}
\toprule
Native theory & Function congruence retained & Clauses & Literals & Nodes & Bits\\
\midrule
FC15$^*$ & All nine (published control) &29&82&320&1616\\
FC15$^*$ & $V,C,N,A,Q$ &25&67&268&1351\\
FC15$^*$ & $V,C$ &22&58&238&1195\\
FC15$^*$ & $V,C$; reconstruct $G$ &21&54&230&1148\\
FC15$^*$ & Also derive reflexivity &20&53&228&1134\\
\midrule
FC16$^*$ & All nine (published control) &30&86&368&1823\\
FC16$^*$ & $V,C,N,A,Q$ &26&71&316&1558\\
FC16$^*$ & $V,C$ &23&62&286&1402\\
FC16$^*$ & $V,C$; reconstruct $G$ &22&58&278&1355\\
FC16$^*$ & Also derive reflexivity &21&57&276&1341\\
\bottomrule
\end{tabular}
\end{center}
All counts include equality completion. Native rows alone remain 15/35
or 16/39 clauses/literals, so mixing a native count with a completed
count would be misleading. The five-base control removes exactly four
witness congruence rows, 15 literals and 265 bits. The two-assignment
target removes seven rows, 24 literals, 82 term nodes and 421 bits.
Its packed FC15$^*$ payload occupies 150 bytes with five trailing zero
padding bits; FC16$^*$ occupies 176 bytes with six padding bits. Maximum
local variable count remains six in the completed templates because $C$
congruence is ternary. Maximum term depths remain three and six.

With both further redundancies removed, the saving is nine completion
rows, 29 literals, 92 term nodes and 482 bits relative to the published
controls. The final packed payloads are 142 bytes with two padding bits
for FC15$^*$ and 168 bytes with three padding bits for FC16$^*$. Every
intermediate bit string is retained as a reproduction control.

The source codec's independent recursive decoder parses the emitted stream
and is compared with the original canonical clause trees. The bit strings,
JSON trees, packed bytes and metrics are saved in
\texttt{skolem\_reduct\_results/}. SHA256 values of the pure bit strings
(without the text file's trailing newline) are
\begin{quote}\small\ttfamily
FC15: 8ffa8bcc38eb03ced7ee019a098988cdc52645112722ada8ceacbbfa3b852169\par
FC16: 6c875817d9ff823b8bd1f6dc6271d0e1f5a88898da9fbabd440bdb0c81e7d0d5.
\end{quote}
For the final \texttt{raw\_guard\_derived\_reflexivity} streams, the
corresponding hashes are
\begin{quote}\small\ttfamily
FC15: 81b2ba36e43c0da53961bb1ca059937e749d0624005876fc921bb7485db59bba\par
FC16: 4a696a464ff72edf35072d8c0eaace39f51a408c27e0ecde5ae9f9e5a0d7758e.
\end{quote}

\section{A matching finite quotient checker}

The existing ProveIt finite-quotient frontend represents a ground valuation
by class labels and three predicate tables. For $k=0,1,\ldots$, put
$B=2^k$. A candidate integer has exactly $kB+3B^2$ bits. Its low field
contains labels $\pi(i)\in\{0,\ldots,B-1\}$ for $i<B$; the rest contains
the $E,\srInc,G$ tables on label pairs. Equality is label equality.
Thus reflexivity, the left-Euclidean law and predicate congruence require
no separately instantiated runtime clauses.

For the reduced target, the sole change to the function check is to retain
tags 0 and 1, $V$ and $C$, instead of all tags 0 through 8. If two in-range
applications of a retained function have coordinatewise equal argument
labels, their output labels must agree. The other seven proper functions
and the seven padding constructors are still in the total ground term
codec; only their congruence obligations are absent. Native terms and all
colliding variable assignments are unchanged.

At stage $k$ retain a native ground instance iff all its endpoint term
codes are below $B$. The radix16/Cantor codec has children below parents,
so the code interval is subterm-closed. Every native ground instance is
eventually retained as $B$ grows. Each candidate is checked against the
retained function pairs and native instances. Exhausting the entire finite
candidate range is UNSAT; a time limit, memory failure or incomplete run
is not an UNSAT certificate.

The same finite-label argument as in the existing frontend proves that a
candidate exists iff the corresponding finite ground CNF with the reduced
completion is satisfiable. Any satisfying valuation labels an object by
the least member of its equality class and stores the induced predicate
tables. Conversely any valid labeling gives all the equality and predicate
rows and the two checked function rows. These finite CNFs are nested, and
every ground clause occurs in some stage. Propositional compactness and
Theorem~\ref{sr-main}, or its FC16$^*$ version, therefore establish in ZFC
\[
\text{every stage has a valid candidate}
\quad\Longleftrightarrow\quad\srCon(\mathrm{ZFC}).
\]
Consequently the ideal program that halts at the first failed stage has
the requested consistency contract. This is an executable specification
to compile, not a finished small binary machine.

The optional \texttt{raw\_guard=True} mode prices the guard corollary
without requiring a larger candidate field. Keep class-indexed lookup for
$E$ and $\srInc$, but use the original endpoint term codes $(a,b)$ as the
indices of the $G$ table. That table still occupies $B^2$ bits. Its lookup
rule is
\[
G(a,b)=\operatorname{bit}\bigl(W,\,kB+2+3(Ba+b)\bigr).
\]
No $G$ congruence is then silently imposed. A two-literal fixture has
$\pi(0)=\pi(1)=0$, $G(0,0)$ true and $G(1,0)$ false: the raw-guard
mode accepts and the original class-indexed mode rejects, as required.

For the reflexivity corollary, finite label certificates still impose
reflexivity structurally. At a small cutoff the displayed two-premise
derivation may contain the larger term $p_z$ outside the interval, so
this must be described as an admissibly strengthened stage sequence.
It need not have the same satisfiability result as the unstrengthened
finite subset at each cutoff. At the all-ground limit, however,
reflexivity is a theorem of the reduced native theory. Adding it at every
finite stage changes neither global ground satisfiability nor the
compactness argument. This proves the same consistency contract for the
final $U_{15}^{\circ},U_{16}^{\circ}$ targets.

The bundled arithmetic function checker was independently compared with
full decoded-term grouping on 1,440 random finite labelings. Local fixtures
separate required from removable congruence: a noncongruent $W$ application
is accepted in the two- and five-function modes and rejected in the
nine-function mode; a noncongruent $N$ application is accepted only in the
two-function mode; a noncongruent $V$ application is rejected. These have
empty native clause lists and are explicitly \emph{local controls}, not
finite models of FC15$^*$. The recorded native stages $k=0,1$ contain no
kept native instances, so their SAT outputs are only staging regressions.
The inherited explicit native $k=2$ witness, $2{,}097{,}728$, is also
checked against twelve retained native clauses in each theory, and the
zero candidate is rejected. This does not exhaust the $2^{56}$-candidate
range or check any FC16$^*$ extensionality instance at this small cutoff.

For a concrete code-interval comparison, the numbers of same-tag pairs
examined by the function condition are
\begin{center}
\begin{tabular}{rrrr}
\toprule
$B$& $V,C$ & $V,C,N,A,Q$ & All nine\\
\midrule
16&0&0&0\\
32&2&5&9\\
64&12&30&54\\
128&56&140&252\\
\bottomrule
\end{tabular}
\end{center}
These counts measure explicit obligations at equal cutoffs, not runtime or
machine states. Template emission, term arithmetic, label extraction,
finite enumeration, register allocation and blank-tape initialization all
remain part of a complete compiler budget. A nine-function signature and
a 482-bit static reduction alone give no numerical upper bound below 245.

\section{Why membership transport remains: a full FC15 countermodel}

The retained $E$-congruence row cannot be removed from FC15$^*$ by the
same argument. There is a concrete model in ZFC of \emph{all fifteen
native rows and all the original completion rows except $E$ transport}.
This strengthens the simple observation that the atom-profile argument
does not by itself transport raw $E$ endpoints.

Take the carrier $\mathbb N\times\{0,1\}$, write $\beta(x)$ for its first
coordinate, and put $\iota(n)=(n,0)$. Let $p_j$ be the $j$th prime,
starting at $p_0=2$, and define integer assignment operations
\[
v(s,j)=\nu_{p_j}(s+1),\qquad
c(s,j,z)=\frac{s+1}{p_j^{v(s,j)}}p_j^z-1.
\]
The quotient in the definition of $c$ is an integer. Changing one prime
exponent gives $v(c(s,j,z),j)=z$ and preserves $v(s,k)$ when $k\ne j$.
Interpret the function symbols by
\begin{align*}
d&=\iota(0),\\
V(s,j)&=\iota(v(\beta(s),\beta(j))),\\
C(s,j,z)&=\iota(c(\beta(s),\beta(j),\beta(z))),\\
N(a,b)&=\begin{cases}\iota(0)&\beta(a),\beta(b)\text{ both odd},\\
                         \iota(1)&\text{otherwise},\end{cases}\\
A(j,k)&=\iota(0),\qquad Q(j,a)=\iota(\beta(a)),\\
W(s,j,a)&=(0,1),\qquad H(b,t)=\iota(0),\\
Y(b,t,u)&=Z(b,t,u)=(\beta(u)+1,1).
\end{align*}
Interpret the predicates by
\begin{align*}
\srEq(x,y)&\quad\Longleftrightarrow\quad\beta(x)=\beta(y),\\
\srInc(s,a)&\quad\Longleftrightarrow\quad\beta(a)\text{ is odd},\\
G(b,t)&\quad\Longleftrightarrow\quad\beta(t)\text{ is even},\\
E((n,i),(m,j))&\quad\Longleftrightarrow\quad n=m\ \land\ i=j=1.
\end{align*}
Every function's first output coordinate depends only on the first input
coordinates. Thus \emph{all nine} function-congruence rows hold, as do
equivalence and the $\srInc,G$ predicate rows.

The prime-exponent laws verify native rows 01--02. NAND on parity verifies
03--05. Every $V$ output has decoration zero, so $E(V(s,j),V(s,k))$ is
false, agreeing with the constant-zero atom constructor in rows 06--07.
Incidence is independent of its assignment argument, and $Q$ preserves
its class-code parity, so rows 08--09 hold for the given constant $W$.

For the closure rows, $R(b,t,x,y)$ is simply ``$\beta(t)$ is odd,'' independent
of $x,y$. When $\beta(t)$ is even, $G$ is true and $R$ is empty. Rows 13--14
are true through $\neg R$, and row 15 is true because $H=\iota(0)$ has
no incoming $E$-member. When $\beta(t)$ is odd, $G$ is false and $R$ is full.
Then rows 13--15 are guarded, row 10 is true through $R$, row 11 has the
true diagonal edge $E(Y,Y)$, and row 12 is true because the first
coordinate of $Z$ is one larger than that of $u$. This exhausts all rows.

The failure of transport is genuinely ground. Let $w=W(d,d,d)=(0,1)$.
Then $\srEq(w,d)$ and $E(w,w)$ hold, while $E(d,d)$ fails. The structure
therefore refutes the $E$-congruence row on these named terms.
Since ZFC constructs this model directly, it proves consistency of the
corresponding weakened finite FC15 theory. Assuming ZFC is consistent,
that weakened theory cannot have a ZFC-provable consistency equivalence
to ZFC, by the second incompleteness theorem.

This obstruction is deliberately scoped. It is \emph{not} a model of
the weakened FC16$^*$ target: $q=\iota(1)$, and the ground selector
$j=C(d,q,q)=\iota(2)$ avoids $d,q$. At $a=d,b=q$, all three equality
guards of the added clause fail and its XOR incidence is false.
No conclusion about removing $E$ congruence specifically from FC16$^*$
is drawn. The delivered positive targets keep that row in both theories.
The independent review includes exact evaluation of the original native
templates on bounded samples of this infinite structure; that finite
diagnostic supports the specification and does not replace the general
proof above.

\section{Reproduction, provenance and proof status}

Run, with the standard Python library only,
\begin{verbatim}
python skolem_reduct.py --output skolem_reduct_results
\end{verbatim}
The generator writes exact payloads and \texttt{results.json}; a failed
assertion raises an error. The top-level checker does not wrap resource
errors as logical conclusions. The quotient argument is an ordinary
mathematical proof with explicit ZFC-internal constructions. It has not
been serialized as a Lean or Coq theorem, and no claim about CPython's
unbounded physical execution is part of the theorem.

This contribution extends the earlier source44 observation that unused
padding functions can be forgotten before quotienting. The additional
step is to reconstruct \emph{used} class and witness functions by their
particular contracts while preserving only assignment operations. It does
not justify deleting equality transport in the existing 245-state
Hilbert/Tarski backend; that is a different proof system, and prior
typed countermodels prohibit several such deletions. No claim of
historical priority beyond the inspected repository is made.

Primary repository sources, pinned for this delivery to ProveIt commit
\texttt{6446ff5668307a2ac6aee1896037f0b792090389}, are:
\begin{enumerate}
\item Source44 mathematical specification and calibration:
\url{https://github.com/VladimirReshetnikov/ProveIt/blob/6446ff5668307a2ac6aee1896037f0b792090389/SetTheory/Cardinals/docs/reports/set-theory/zfc-closure-kernels/44-nine-notes-finite_class_closure.md}.
\item Exact native clause and equality-completion generator:
\url{https://github.com/VladimirReshetnikov/ProveIt/blob/6446ff5668307a2ac6aee1896037f0b792090389/SetTheory/Cardinals/docs/reports/set-theory/zfc-closure-kernels/code/44-nine-code-ground_kernel.py}.
\item Exact FC16$^*$ extension and independent prefix decoder:
\url{https://github.com/VladimirReshetnikov/ProveIt/blob/6446ff5668307a2ac6aee1896037f0b792090389/SetTheory/Cardinals/docs/reports/set-theory/zfc-closure-kernels/code/44-nine-code-coded_extensional.py}.
\item Accepted source audit, including inherited calibration boundaries:
\url{https://github.com/VladimirReshetnikov/ProveIt/blob/6446ff5668307a2ac6aee1896037f0b792090389/SetTheory/Cardinals/docs/reports/set-theory/zfc-closure-kernels/intake-155-source-44.md}.
\item Existing finite quotient representation and checker:
\url{https://github.com/VladimirReshetnikov/ProveIt/blob/6446ff5668307a2ac6aee1896037f0b792090389/Computability/BusyBeaver/Research/ZFC/finite_quotient_frontend.md}.
\end{enumerate}
The general finite assignment/class framework is attributed by source44
to Harvey Friedman's \emph{Impossible Counting}, especially its
Section~4.1; the present source44 native clauses and set-theoretic closure
synthesis retain their original attribution. Byte-exact copies of the
two generator modules, the source44 mathematical note and the repository
license are in \texttt{reference/}; their hashes are included in the
generated result.
